% GENERATED_BY: run_oc_core_monograph_translation_rework_dispatch_v1.py
% AI_ASSISTED_TRANSLATION_DRAFT: TRUE
% THEOREM_REVIEW_STATUS: PENDING
% THEOREM_REVIEW_MODE: FORMAL_AI_GATE__CODEX_CLI_CHATGPT
% THEOREM_REVIEW_REVIEWER_ID: 
% THEOREM_REVIEW_AUTHORITY_CLASS: 
% THEOREM_REVIEWED_AT: 
% THEOREM_REVIEW_DOSSIER_REF: 
% SOURCE_LANGUAGE: EN
% TARGET_LANGUAGE: DE
% SOURCE_EN_REF: content/04_results.tex
\section{Ergebnisse und abgeleitete Konsequenzen}
\label{sec:results}

Dieser Abschnitt fasst die zentralen strukturellen Resultate zusammen, die aus dem in Abschnitt~\ref{sec:model} eingeführten formalen Rahmen folgen.
Anders als in früheren Entwürfen des Kerns ist dieses Kapitel kein Platzhalter: Es präsentiert die kanonischen Konsequenzen der Axiome und Operatoren der Ontology of Continua (OC), konsolidiert aus früheren internen Durchläufen und für Kern~1.3 neu geordnet.

Alle Resultate dieses Kapitels sind domänenunabhängig.
Sie folgen ausschließlich aus den strukturellen Definitionen von Kontinua, Achsen, Schwellen, Potenzialen, Flüssen, Grenzen, Zyklen und dem Maß der Kontinuumshaftigkeit \(k(K,t)\).
Die Beweise werden hier nicht vollständig wiedergegeben; sie sind eigenständigen Erweiterungsarbeiten vorbehalten.
Im Mittelpunkt steht die Formulierung der zentralen Sätze und Konsequenzen, die alle Kontinua in der Hierarchie \(K_0\)–\(K_{10}\) beschränken.

\subsection{Satz 1: Monotonie der Dimensionalität}

\textbf{Aussage.}
Wenn ein Kontinuum \(K_{x+1}\) aus \(K_x\) hervorgeht, dann
\[
    \dim(K_{x+1}) > \dim(K_x).
\]

\textbf{Begründung (Skizze).}
Dimensionale Emergenz ist als das Auftreten einer neuen Achse \(A_{\mathrm{new}}\) definiert, die mit allen bestehenden Achsen in \(A(K_x)\) inkompatibel ist; das heißt, sie kann nicht als lineare oder strukturelle Kombination aus ihnen dargestellt werden.
Die dimensionale Schwelle \(\Theta_{\mathrm{dim}}(K_x)\) erzwingt, dass \(A_{\mathrm{new}}\) nicht entartet und strukturell notwendig ist.
Daher hat jedes emergente Kontinuum eine strikt höhere Dimensionalität als sein Vorgänger.

\textbf{Korollar 1.1 (Keine degradierende Vereinfachung).}
Es gibt keine Entwicklung auf Kontinuumsebene, die unter Erhaltung der Existenz die Dimensionalität verringert:
ist \(K(t)\) lebendig (\(k(K,t) > 0\)), dann
\[
    \dim\big(K(t+dt)\big) \ge \dim\big(K(t)\big).
\]
Eine scheinbare Verringerung der Dimension entspricht dem Tod des höherdimensionalen Kontinuums und der Geburt eines anderen, niedrigerdimensionalen Kontinuums, nicht einer kontinuierlichen Vereinfachung.

\subsection{Satz 2: Unmöglichkeit spontaner Dimensionsentstehung}

\textbf{Aussage.}
Kein Kontinuum kann seine Dimensionalität erhöhen, ohne die dimensionale Schwelle zu überschreiten und Zugang zu einer geeigneten Achse in seinem Einbettungsraum zu haben:
\[
    \dim(K_{x+1}) > \dim(K_x)
    \quad\Rightarrow\quad
    T(K_x,t) > \Theta_{\mathrm{dim}}(K_x)
    \ \wedge\
    A_{\mathrm{new}} \in A(M_x)\setminus A(K_x).
\]

\textbf{Begründung (Skizze).}
Neue Achsen repräsentieren Klassen von Unterschieden, die mit vorhandenen Achsen inkompatibel sind.
Eine solche Inkompatibilität entsteht nur dann, wenn die strukturelle Spannung \(T(K_x,t)\) im Verhältnis zur aktuellen Schwellenlandschaft nicht innerhalb der vorhandenen Achsen aufgelöst werden kann.
Nach Definition von \(\Theta_{\mathrm{dim}}\) können unterhalb dieser Schwelle alle Unterschiede auf den Spann von \(A(K_x)\) projiziert werden; oberhalb davon scheitert die Projektion, und eine neue Achse ist erforderlich.
Die neue Achse muss im Einbettungsraum \(M_x\) bereits verfügbar sein; andernfalls gibt es keine strukturelle Richtung, entlang derer sich das Kontinuum ausdehnen kann.

\textbf{Folgerung 2.1.}
Rauschen, lokale Fluktuationen oder interne Flüsse, die die strukturelle Spannung nicht über \(\Theta_{\mathrm{dim}}(K_x)\) anheben, können keine neuen Dimensionen erzeugen.

\subsection{Satz 3: Tod durch Grenz- und Einbettungskollaps}

\textbf{Aussage.}
Ein Kontinuum \(K\) stirbt, wenn sein zulässiger Zustandsraum kollabiert:
\[
    \Omega(K) = \emptyset.
\]

\textbf{Äquivalente Formulierungen.}
Die folgenden Bedingungen sind äquivalent und charakterisieren den Tod von \(K\):
\begin{enumerate}
    \item tragende Zyklen verschwinden, \(C(K) = \emptyset\);
    \item tragende Flüsse verschwinden, \(J_{\mathrm{support}} = 0\), oder reichen nicht aus, um irgendwelche Zyklen aufrechtzuerhalten;
    \item es gibt eine zustandsunabhängige Verletzung von Schwellen, \(\forall s:\ \exists\,k\ \text{mit } f_k(s) > 0\);
    \item die Kontinuumshaftigkeit kollabiert, \(k(K,t) \to 0\);
    \item das Kontinuum verlässt den von seinem Einbettungsraum \(M\) getragenen Bereich, d.\,h.\ kein Zustand erfüllt zugleich die internen Schwellen von \(K\) und die externen Beschränkungen von \(M\).
\end{enumerate}

\textbf{Korollar 3.1 (Tod als Austritt aus dem Einbettungsraum).}
Wenn sich die Beschränkungen von \(M\) so ändern, dass keine Konfiguration von \(K\) unter Einhaltung seiner Schwellen in \(M\) eingebettet werden kann, dann wird \(\Omega(K)\) leer und das Kontinuum stirbt, unabhängig von seiner internen Struktur.

\textbf{Korollar 3.2 (Irreversibilität des Todes).}
Sobald \(\Omega(K) = \emptyset\), kann weder ein Operator \(E\), der innerhalb derselben Ebene wirkt, noch irgendein Interaktionsoperator \(E_{\mathrm{int}}\), der die Identität von \(K\) bewahrt, ein nichtleeres \(\Omega(K)\) rekonstruieren.
Jede scheinbare ``Wiederherstellung'' entspricht der Geburt eines neuen Kontinuums \(K'\), nicht der Auferstehung des ursprünglichen \(K\).

\subsection{Satz 4: Unmöglichkeit der Evolution außerhalb des Einbettungsraums}

\textbf{Aussage.}
Sei \(K_x\) ein in \(M_x\) eingebettetes Kontinuum.
Dann gilt für alle Zeiten \(t\),
\[
    A\big(K_x(t)\big) \subseteq A(M_x),
    \qquad
    \mathrm{span}\,A\big(K_x(t+dt)\big) \subseteq \mathrm{span}\,A(M_x).
\]

\textbf{Folgerung 4.1.}
Alle dynamischen Prozesse von \(K_x\), einschließlich der dimensionalen Geburt \(K_x \to K_{x+1}\), erfordern, dass der Einbettungsraum \(M_x\) die Achsen, entlang derer die Entwicklung verläuft, bereits enthält.
Ein Kontinuum kann sich nicht in Richtungen entwickeln, die in seinem Einbettungsraum nicht vorhanden sind.

\subsection{Satz 5: Notwendigkeit und Hinreichendheit der Kompatibilität mit \texorpdfstring{$M$}{M}}

\textbf{Aussage.}
Ein Kontinuum \(K\) existiert als lebendiges Kontinuum (d.\,h.\ \(\Omega(K)\neq\emptyset\) und \(k(K,t)>0\)) genau dann, wenn es mit seinem Einbettungsraum \(M\) kompatibel ist:
\[
    K \text{ existiert }
    \Longleftrightarrow
    \Omega(K) \neq \emptyset
    \ \wedge\
    A(K) \subseteq A(M)
    \ \wedge\
    \Theta(K) \text{ ist innerhalb von } M \text{ erfüllbar}.
\]

\textbf{Folgerung 5.1.}
Selbst wenn eine Konfiguration auf der Ebene von \(K\) mathematisch wohldefiniert ist, entspricht sie keinem lebendigen Kontinuum, sofern der Einbettungsraum die erforderlichen Achsen, Schwellen und Flüsse nicht tragen kann.
Inkompatible Konfigurationen werden strukturell ausgeschlossen, indem \(\Omega(K) = \emptyset\) gesetzt wird.

\subsection{Satz 6: Strukturelle Spannung und Phasenübergänge}

\textbf{Aussage.}
Sei \(T(K,t)\) die strukturelle Spannung, verstanden als ein Funktional der Potenziale, Achsen und ihrer Gradienten:
\[
    T(K,t) = T\big(P(t),A,\nabla P(t)\big).
\]
Dann gilt:
\begin{itemize}
    \item ein Phasenübergang tritt ein, wenn
          \(
              T(K,t) = \Theta_{\mathrm{crit}}(K),
          \)
    \item ein dimensionaler Übergang tritt ein, wenn
          \(
              T(K,t) = \Theta_{\mathrm{dim}}(K),
          \)
    \item ein struktureller Kollaps tritt ein, wenn
          \(
              T(K,t) = \Theta_{\mathrm{death}}(K).
          \)
\end{itemize}

\textbf{Korollar 6.1.}
Alle qualitativen Veränderungen im Kontinuum --- das Auftreten neuer Regime, neuer Dimensionen oder ein Kollaps --- werden durch die Beziehung zwischen struktureller Spannung und der Schwellenlandschaft bestimmt.

\subsection{Satz 7: Universelles Gesetz des Komplexitätswachstums}

\textbf{Aussage.}
Für jedes lebendige Kontinuum \(K\),
\[
    S\big(K(t+dt)\big) \ge S\big(K(t)\big),
\]
wobei \(S\) ein strukturelles Maß der Komplexität ist, das auf dem Zustandsraum und seiner Organisation definiert ist (beispielsweise durch Kombination von Beiträgen aus der Anzahl der Achsen, der Diversität der Zustände, der Reichhaltigkeit der Zyklen und Eigenschaften der Einbettung).

\textbf{Folgerung 7.1.}
Die Komplexität nimmt strikt zu, sobald:
\begin{itemize}
    \item neue Achsen auftreten;
    \item sich neue stabile Zyklen bilden;
    \item sich der zulässige Zustandsraum \(\Omega(K)\) erweitert;
    \item der Einbettungsraum \(M\) neue, für \(K\) relevante Achsen gewinnt.
\end{itemize}
Stagnierende oder konstante Komplexität entspricht fein austarierten Dynamiken unterhalb kritischer Schwellen; dieses Regime ist jedoch im Allgemeinen instabil (siehe Satz~\ref{thm:no-eternal-stabilisation}).

\subsection{Satz 8: Unmöglichkeit ewiger Stabilisierung}
\label{thm:no-eternal-stabilisation}

\textbf{Aussage.}
Es gibt kein nichttriviales lebendiges Kontinuum \(K\), das für immer an einem Fixpunkt seines Zustandsraums verbleibt und dabei positive Kontinuumshaftigkeit aufrechterhält.
Formal existiert keine Lösung mit
\[
    \frac{dP}{dt} = 0, \quad \frac{dJ}{dt} = 0, \quad \frac{d\Theta}{dt} = 0, \quad \frac{dk}{dt} = 0
\]
für alle \(t\), außer im trivialen Fall, dass \(K\) strukturell eingefroren und von seinem Einbettungsraum entkoppelt ist.

\textbf{Begründung (Skizze).}
Einbettungsräume \(M_x\) entwickeln und erweitern sich selbst (Satz~10).
Interne und externe Flüsse können nicht auf allen Zeitskalen zugleich identisch null sein, ohne entweder:
\begin{itemize}
    \item das Kontinuum in den Kollaps zu treiben (Verlust tragender Flüsse), oder
    \item strukturelle Änderungen (Achsen, Schwellen, Zyklen) hervorzurufen.
\end{itemize}
Folglich unterliegt jedes nichttriviale lebendige Kontinuum entweder einer Entwicklung seiner Struktur oder letztlich dem Tod.

\subsection{Satz 9: Tod als Verlust von Zyklen}

\textbf{Aussage.}
Ein lebendiges Kontinuum \(K\) stirbt, wenn sein maximaler strukturell stabiler Zykluskomplex verschwindet:
\[
    C_{\max}(K) = \emptyset.
\]

\textbf{Begründung (Skizze).}
Zyklen repräsentieren selbsterhaltende Flüsse, die das Kontinuum von seiner Grenze \(\partial\Omega\) fernhalten.
Wenn kein solcher Zyklus existiert, dann:
\begin{itemize}
    \item können tragende Flüsse keine geschlossenen Schleifen bilden;
    \item nähern sich Trajektorien entweder \(\partial\Omega\) an oder verletzen Schwellen;
    \item fällt die Kontinuumshaftigkeit \(k(K,t)\) gegen null.
\end{itemize}
Folglich fällt der Verlust des maximalen Zykluskomplexes mit dem Kollaps von \(\Omega(K)\) zusammen.

\textbf{Korollar 9.1.}
Tod kann strukturell durch das Verschwinden aller Zyklen erkannt werden, die eine minimale Stabilitätsbedingung erfüllen (streng positiver Abstand zu \(\partial\Omega\)).

\subsection{Satz 10: Monotones Wachstum von Einbettungsräumen}

\textbf{Aussage.}
Einbettungsräume bilden eine monotone Folge:
\[
    M_0 \subset M_1 \subset \dots \subset M_x \subset \dots,
\]
und immer dann, wenn ein neues Kontinuum \(K_{x+1}\) entsteht, erweitert der entsprechende Einbettungsraum \(M_{x+1}\) den Raum \(M_x\) bezüglich seiner verfügbaren Achsen strikt:
\[
    A(M_{x}) \subset A(M_{x+1}).
\]

\textbf{Folgerung 10.1.}
Das Wachstum von Kontinua in Dimension und Komplexität ist vom Wachstum ihrer Einbettungsräume untrennbar.
Neue Kontinua können nicht ohne eine Erweiterung der in \(M\) kodierten strukturellen Möglichkeiten erscheinen.

\subsection{Satz 11: Inkompatibilität zwischen Schwellen und Ausdrucksfähigkeit}

\textbf{Aussage.}
Ein Kontinuum \(K\) kollabiert, wenn seine Achsen nicht mehr ausreichen, um die erforderlichen Unterschiede auszudrücken:
\[
    \dim\big(A(K)\big) < \dim\big(\mathrm{Differences}(K)\big),
\]
wobei \(\mathrm{Differences}(K)\) die effektive Dimension der strukturell relevanten Unterscheidungen ist, die durch den Einbettungsraum und die Schwellen auferlegt werden.

\textbf{Folgerung 11.1.}
Ein Kollaps kann selbst ohne Verletzung energetischer oder dynamischer Beschränkungen eintreten, wenn die Repräsentationskapazität überschritten wird.
Insbesondere können kognitive, soziale oder metatheoretische Kontinua allein aufgrund des Verlusts expressiver Adäquatheit ihrer Achsen sterben.

\subsection{Satz 12: Unvollständigkeit von Einbettungsräumen}

\textbf{Aussage.}
Für jeden endlichen Einbettungsraum \(M_x\) gibt es potenzielle Kontinua \(K'\), die innerhalb von \(M_x\) nicht realisiert werden können, weil ihre Achsen oder Schwellen mit \(A(M_x)\) oder mit den Beschränkungen von \(M_x\) inkompatibel sind.

\textbf{Folgerung 12.1.}
Kein einzelner Einbettungsraum \(M_x\) ist universell für alle möglichen Kontinua.
Die Folge der Einbettungsräume muss sich selbst erweitern und diversifizieren, um neue strukturelle Regime zu beherbergen.

\subsection{Korollare und domänenunabhängige Konsequenzen}

Aus der Zusammenstellung der obigen Sätze ergeben sich die folgenden domänenunabhängigen Konsequenzen:

\begin{itemize}
    \item Die Dimensionalität ist für lebendige Kontinua strikt monoton; eine kontinuierliche Degradation der Dimension ist nicht möglich.
    \item Dimension kann nicht spontan auftreten; sie erfordert strukturelle Spannung oberhalb von \(\Theta_{\mathrm{dim}}\) und das Vorhandensein geeigneter Achsen im Einbettungsraum.
    \item Tod ist durch den Kollaps von \(\Omega(K)\), das Verschwinden von Zyklen und die Irreversibilität dieses Kollapses gekennzeichnet.
    \item Evolution ist durch Einbettungsräume beschränkt; Kontinua können sich nicht in Richtungen entwickeln, die \(M\) nicht trägt.
    \item Komplexität neigt bei lebendigen Kontinua zum Wachstum, angetrieben durch neue Achsen, neue Zyklen und sich erweiternde Zustandsräume.
    \item Ewige exakte Stabilisierung ist für nichttriviale Kontinua strukturell ausgeschlossen; sie entwickeln sich entweder oder sterben.
    \item Repräsentationale und expressive Grenzen können einen Kollaps verursachen, selbst wenn energetische und dynamische Grenzen noch nicht erreicht sind.
\end{itemize}

Diese Konsequenzen gelten einheitlich für alle Ebenen \(K_0\)–\(K_{10}\); ihr domänenspezifischer Gehalt ergibt sich erst aus der Interpretation von Achsen, Potenzialen, Schwellen und Flüssen in jeder Ebene.

\subsection{Implikationen für die vertikale Hierarchie}

In der vertikalen Hierarchie \(K_0\)–\(K_{10}\) manifestieren sich die strukturellen Sätze wie folgt:

\begin{itemize}
    \item \(K_0\): keine Dynamik, keine Geburt und kein Tod; nur strukturelle Bedingungen für Unterscheidbarkeit.
    \item \(K_1\): klassisches phasenartiges Verhalten mit kontinuierlichen Flüssen und grundlegenden Stabilitäts- und kritischen Schwellen.
    \item \(K_2\): physikalische Schwellen, einschließlich Perkolation, BKT-Übergängen, Kohärenz-/Kondensationsschwellen und massenerzeugenden Mechanismen.
    \item \(K_3\)–\(K_4\): chemische und präbiotische Schwellen (RAF-Abschluss, Membranbildung, osmotische und Krümmungsschwellen, Redox- und pH-Schwellen).
    \item \(K_5\): elektrische Schwellen für Erregung, Spike-Bildung und gradientgetriebene Flüsse in frühen neuronalen Systemen.
    \item \(K_6\): logische und repräsentationale Schwellen für kognitive Bindung, Vorhersage, Gedächtnisstabilität und Kohärenz interner Modelle.
    \item \(K_7\)–\(K_8\): soziale und zivilisatorische Schwellen, die Vertrauen, institutionelle Stabilität, systemische Resilienz und Kollaps steuern.
    \item \(K_9\)–\(K_{10}\): Ausdrucks-, Kohärenz- und Konsistenzschwellen im Raum von Theorien, Paradigmen, Ontologien und metatheoretischen Strukturen.
\end{itemize}

In jedem Fall gelten dieselben strukturellen Prinzipien --- monotone Dimension, schwellengeleitete Emergenz und Kollaps, Abhängigkeit von Einbettungsräumen und Komplexitätswachstum ---, während sich die konkrete Interpretation der Variablen ändert.

\subsection{Zusammenfassung}

Die in diesem Kapitel dargestellten Ergebnisse bilden das strukturelle Rückgrat der OC.
Sie besagen, dass die Dimension monoton ist, dass emergente Dimensionen schwelleninduzierte Phasenübergänge erfordern, dass der Tod durch den Kollaps zulässiger Zustände gekennzeichnet und irreversibel ist, dass Evolution durch Einbettungsräume beschränkt ist, dass Komplexität bei lebendigen Kontinua zum Wachstum neigt, dass Zyklen die Träger strukturellen Lebens sind und dass expressive Grenzen einen Kollaps verursachen können.
Diese Prinzipien gelten über die gesamte Kontinuumshierarchie hinweg und definieren die Beschränkungen, unter denen alle Kontinua existieren, sich entwickeln oder sterben können.
